% Lösungen Einheit 2 von 4 – Äquivalenzumformungen (zusammenbau v0.1)
\einheitenkopf{Einheit 2 von 4 · Äquivalenzumformungen}

\erg{12}{a) $-5$}
\erg{13}{a) $\mid -7$}
%% TODO Lösung der Erklärzeile 14g) fehlt in der Bank
\erg{14}{a) $\mid -11$ \quad b) $x = 7$ \quad c) $x = 17$ \quad d) $x = 9$ \quad e) $x = 14$ \quad f) $x = 37$ \quad g) \ldots \quad h) $x = 9$ \quad i) $x = -7$ \quad j) $x = 7$; Probe: $3 \cdot 7 + 5 = 26$ (wA) \quad k) $x = -6$ \quad l) $x = 24$ \quad m) zum Beispiel $2x + 3 = 15$ (Probe: $2 \cdot 6 + 3 = 15$) \quad n) $x + 7 = 4$, $x = -3$; Probe: $3 \cdot (-3 + 7) = 12$ (wA)}
\erg{15}{a) Bei $3x$ wird mit 3 malgenommen, die Umkehrung ist $:3$, nicht $-3$. Richtig: $x = 5$. \quad b) Die 8 muss auf beiden Seiten abgezogen werden: $x = 14 - 8 = 6$. \quad c) Beim Teilen durch 3 wird jedes Glied geteilt: $x + 3 = 7$, also $x = 4$ (einfacher: erst $-9$, dann $:3$). \quad d) Beide Seiten sind gleich groß; zieht man von beiden dasselbe ab, bleiben sie gleich groß – die Lösung ändert sich nicht.}
